% ============================================================
%  Chapter 4 — The Conical Helix: Direct Construction
% ============================================================
\section{The Conical Helix}\label{sec:helix}

Using the universal quantity $Q(k) = \sqrt{48k+1}$ from
Section~\ref{sec:identity}, we now construct the conical helix
directly and prove its geometric properties as algebraic corollaries.

\subsection{Exact Parametrisation}

\begin{theorem}[Exact helix parametrisation]\label{thm:helix}
Let $Q(k) = \sqrt{48k+1}$. The conical helix is given exactly by
\begin{equation}\label{eq:helix-exact}
  \boxed{
  \vec{p}(k) =
  \begin{pmatrix}
    \dfrac{Q}{2\pi}\cos\!\Bigl(\dfrac{\pi\,Q}{12}\Bigr)\\[10pt]
    \dfrac{Q}{2\pi}\sin\!\Bigl(\dfrac{\pi\,Q}{12}\Bigr)\\[10pt]
    \dfrac{Q-3}{4}
  \end{pmatrix},
  \qquad Q = Q(k).
  }
\end{equation}
At the integer values $k = k(n)$ with $n \equiv 1, 5 \pmod{6}$,
one has $Q = 4n+3$, and the point $\vec{p}$ has integer height $z = n$.
\end{theorem}

\begin{proof}
By Definition~\ref{def:Q} and Table~\ref{tab:Q-overview}, the three
coordinates are $r = Q/(2\pi)$, $\theta = \pi Q/12$, and $z = (Q-3)/4$.
At $k = k(n)$ the fundamental identity (Theorem~\ref{thm:fundamental})
gives $Q = 4n+3$, whence $z = (4n+3-3)/4 = n$ and
$\theta = \pi(4n+3)/12 = \pi n/3 + \pi/4$.
\end{proof}

Equivalently, the helix admits the standard $\theta$-parametrisation:

\begin{corollary}[$\theta$-parametrisation]\label{cor:theta-param}
With $a = 6/\pi^2$ and $c_1 = 3/\pi$:
\begin{equation}\label{eq:helix-theta}
  \gamma(\theta) =
  \begin{pmatrix}
    a\,\theta\cos\theta \\[4pt]
    a\,\theta\sin\theta \\[4pt]
    -\tfrac{3}{4} + c_1\,\theta
  \end{pmatrix}.
\end{equation}
The helix intersects the plane $z = n$ at the angle
$\theta_n = \frac{\pi}{3}n + \frac{\pi}{4}$.
\end{corollary}

\begin{proof}
Setting $\theta = \pi Q/12$ and $r = Q/(2\pi)$ gives
$r = (6/\pi^2)\theta = a\theta$. The height is
$z = (Q-3)/4 = (12\theta/\pi - 3)/4 = (3/\pi)\theta - 3/4 = c_1\theta - 3/4$.
\end{proof}

\subsection{Spiral Conformity}\label{subsec:spiral}

\begin{theorem}[Archimedean spiral]\label{thm:spiral}
The helix points $(r(k), \theta(k))$ lie exactly on the Archimedean
spiral $r = \frac{6}{\pi^2}\,\theta$.
\end{theorem}

\begin{proof}
Direct verification:
$\frac{6}{\pi^2}\cdot\frac{\pi Q}{12}
  = \frac{Q}{2\pi}
  = r(k)$.
\end{proof}

\begin{remark}[Connection to the Basel problem]
The spiral slope constant $a = 6/\pi^2 = 1/\zeta(2)$ is the reciprocal
of the Riemann zeta function at $s = 2$, the solution of Euler's
Basel problem~\cite{Euler1740, Hardy2008}.
\end{remark}

\subsection{Cone Conformity}\label{subsec:cone}

\begin{theorem}[Cone equation]\label{thm:cone}
The helix points $(r(k), z(k))$ lie exactly on the cone of revolution
$z = -\frac{3}{4} + \frac{\pi}{2}\,r$.
\end{theorem}

\begin{proof}
Direct verification:
$-\frac{3}{4} + \frac{\pi}{2}\cdot\frac{Q}{2\pi}
  = -\frac{3}{4} + \frac{Q}{4}
  = \frac{Q-3}{4}
  = z(k)$.
\end{proof}

\begin{remark}[Geometric interpretation]
The cone has apex at $(0, 0, -3/4)$, half-angle
$\alpha = \arctan(2/\pi) \approx 32.48°$, and the helix rises by
exactly $\Delta z = 6$ per full revolution
($\Delta\theta = 2\pi$, $\Delta z = c_1 \cdot 2\pi = 6$).
\end{remark}

\subsection{Three Invariants from One Identity}

\begin{tcolorbox}[
  title={Summary: Three simultaneous invariants},
  colback=blue!3,
  colframe=blue!40!black,
  enhanced
]
The parametrisation~\eqref{eq:helix-exact} satisfies three exact
invariants, all algebraic consequences of the single identity
$(4n+3)^2 = 48k(n)+1$:
\begin{enumerate}[label=(\roman*)]
  \item \textbf{Arithmetic invariant:} $Q(k(n)) = 4n+3$
        \quad (Theorem~\ref{thm:fundamental}),
  \item \textbf{Spiral invariant:} $r = \frac{6}{\pi^2}\,\theta$
        \quad (Theorem~\ref{thm:spiral}),
  \item \textbf{Cone invariant:} $z = -\frac{3}{4} + \frac{\pi}{2}\,r$
        \quad (Theorem~\ref{thm:cone}).
\end{enumerate}
\end{tcolorbox}

\subsection{Velocity, Curvature, and Torsion}\label{subsec:curvature}

The speed of the helix in the $\theta$-parametrisation is obtained by
differentiating~\eqref{eq:helix-theta}. The cross terms cancel via the
Pythagorean identity, yielding:
\begin{equation}\label{eq:speed}
  \|\gamma'(\theta)\|
  = \frac{3}{\pi^2}\sqrt{4(1+\theta^2) + \pi^2}.
\end{equation}

\begin{theorem}[Exact 3D curvature and torsion]\label{thm:curvature-torsion}
With $a = 6/\pi^2$ and $c_1 = 3/\pi$:
\begin{align}
  \kappa(\theta)
  &= \frac{a\,\sqrt{c_1^2(\theta^2+4) + a^2(\theta^2+2)^2}}
          {\bigl(a^2(1+\theta^2) + c_1^2\bigr)^{3/2}},
  \label{eq:kappa}\\[6pt]
  \tau(\theta)
  &= \frac{c_1\,(\theta^2 + 6)}
          {c_1^2(\theta^2+4) + a^2(\theta^2+2)^2}.
  \label{eq:tau}
\end{align}
\end{theorem}

\begin{proof}
Computing the derivatives $\gamma'$, $\gamma''$, $\gamma'''$ of the
parametrisation~\eqref{eq:helix-theta}, one obtains the cross product
\[
  \gamma' \times \gamma''
  = \bigl(-ac_1(2\cos\theta - \theta\sin\theta),\;
          -ac_1(2\sin\theta + \theta\cos\theta),\;
          a^2(\theta^2+2)\bigr)
\]
with $|\gamma' \times \gamma''|^2 = a^2[c_1^2(\theta^2+4) +
a^2(\theta^2+2)^2]$, $|\gamma'|^2 = a^2(1+\theta^2) + c_1^2$, and
the scalar triple product
$(\gamma' \times \gamma'') \cdot \gamma''' = a^2 c_1(\theta^2+6)$.
The formulae follow from $\kappa = |\gamma' \times \gamma''|/|\gamma'|^3$
and $\tau = (\gamma' \times \gamma'') \cdot \gamma''' / |\gamma' \times
\gamma''|^2$.
\end{proof}

\begin{remark}[Asymptotics]
For $\theta \to \infty$:
$\kappa(\theta) \sim 1/(a\theta) \to 0$ (the helix straightens) and
$\tau(\theta) \to c_1/a^2 = \pi^3/12$ (the torsion converges to a
non-zero constant).
\end{remark}

\subsection{Angular Spacing Pattern}\label{subsec:angles}

\begin{proposition}[Alternating angular spacings]\label{prop:angles}
Between consecutive integer points the angular spacings are:
\begin{align}
  \Delta\theta^{(4)} &= \theta_{n+4} - \theta_n = \tfrac{4\pi}{3}
  \quad\text{(long step)},
  \label{eq:angle-long}\\
  \Delta\theta^{(2)} &= \theta_{n+2} - \theta_n = \tfrac{2\pi}{3}
  \quad\text{(short step)}.
  \label{eq:angle-short}
\end{align}
Together, $\Delta\theta^{(4)} + \Delta\theta^{(2)} = 2\pi$,
so that each long--short pair completes exactly one full revolution.
\end{proposition}

\begin{proof}
From $\theta_n = \frac{\pi}{3}n + \frac{\pi}{4}$:
$\theta_{n+\Delta n} - \theta_n = \frac{\pi}{3}\Delta n$.
\end{proof}
